\chapter{Lecture 14: System Modelling, UML and Inheritance}
\begin{center}
\textit{February 27, 2018}
\end{center}

\section{System Modelling}
It is a good idea to sketch out a design before you code. Designing a program
means
\begin{itemize}
  \item identifying the key entities (classes), and
  \item identifying the relationships between classes --- this is the most
        important part.
\end{itemize}

\section{UML}
\textbf{UML} (Unified Modelling Language) is a standard notation for drawing
class diagrams. A class is drawn as a box with three sections: name, fields,
methods. Visibility is marked with \code{+} (public), \code{-} (private),
\code{\#} (protected).

\begin{center}
\begin{tikzpicture}
\node[draw, rectangle split, rectangle split parts=3, text width=4.2cm,
  align=left, font=\small] {
  \centering\textbf{Vec}
  \nodepart{second} \code{-x: Integer}\\ \code{-y: Integer}
  \nodepart{third} \code{+getX(): Integer}\\ \code{+getY(): Integer}};
\end{tikzpicture}
\end{center}

Fields and methods may be omitted when they are not relevant to the diagram.

\section{Relationship 1: Composition (owns-a)}
\begin{lstlisting}[style=cpp]
class Vec {
    int x, y;
  public:
    Vec(int x, int y) : x{x}, y{y} {}
};

class Basis {
    Vec v1, v2;
};

Basis b;   // won't compile: Vec has no default ctor
\end{lstlisting}
To make it compile, either
\begin{enumerate}
  \item give \code{Vec} a default ctor, or
  \item have \code{Basis}'s ctor use the MIL to call a non-default Vec ctor:
\begin{lstlisting}[style=cpp]
class Basis {
    Vec v1, v2;
  public:
    Basis() : v1{1, 0}, v2{0, 1} {}
};
Basis b;   // compiles
\end{lstlisting}
\end{enumerate}

Embedding one object inside another is called \textbf{composition}: a Basis
\textbf{owns-a} Vec (two, in fact). If A owns-a B, then typically
\begin{enumerate}
  \item B has no identity outside A;
  \item if A is destroyed, B is destroyed;
  \item if A is copied, B is copied (deep copy).
\end{enumerate}
Example: a car owns its wheels. In UML, composition is drawn with a
\emph{filled} diamond at the owner's end; multiplicities go at the other end.

\begin{center}
\begin{tikzpicture}[cls/.style={draw, minimum width=20mm, minimum height=8mm}]
\node[cls] (car) {Car};
\node[cls, right=35mm of car] (w) {Wheel};
\draw[-{Diamond[fill=black,length=4mm,width=2.5mm]}] (w) -- node[above, pos=0.15] {4} (car);
\end{tikzpicture}
\end{center}

\section{Relationship 2: Aggregation (has-a)}
Example: parts in a catalogue, ducks in a pond. A catalogue \textbf{has-a}
collection of parts. If A has-a B, then typically
\begin{enumerate}
  \item B exists independently of A;
  \item destroying A does not destroy B;
  \item copying A does not copy B (shallow copy: the copies share the B's).
\end{enumerate}
\begin{lstlisting}[style=cpp]
class Catalogue {
    Part *parts[100];   // pointers to parts that live elsewhere
    ...                 // ~Catalogue() does NOT delete them
};
\end{lstlisting}
In UML, aggregation is drawn with a \emph{hollow} diamond.
\begin{center}
\begin{tikzpicture}[cls/.style={draw, minimum width=20mm, minimum height=8mm}]
\node[cls] (pond) {Pond};
\node[cls, right=35mm of pond] (d) {Duck};
\draw[-{Diamond[open,length=4mm,width=2.5mm]}] (d) -- node[above, pos=0.15] {0..*} (pond);
\end{tikzpicture}
\end{center}

Whether a pointer field means ownership or not is not visible in the code; it is
a design decision that must be documented (and reflected in the Big Five).

\section{Relationship 3: Inheritance (is-a)}
Suppose we have classes \code{Book} (title, author, numPages), \code{Text}
(additionally a topic), and \code{Comic} (additionally a hero). As three
unrelated classes they capture no relationship.

\paragraph{Q.} How do we create an array that can hold all of these?

In C we would use an array of \code{void*}, or an array of a \code{union} type:
\begin{lstlisting}[style=cpp]
union BookTypes {
    Book *b;
    Text *t;
    Comic *c;
};
BookTypes myBooks[100];
\end{lstlisting}
Neither is type-safe: nothing records which kind of book each element is.

Observe: a Text \emph{is a} Book (with a topic field), and a Comic \emph{is a}
Book (with a hero field). Model this with inheritance:
\begin{lstlisting}[style=cpp]
class Book {                       // base class / superclass
    string title, author;
    int numPages;
  public:
    Book(const string &t, const string &a, int n)
        : title{t}, author{a}, numPages{n} {}
};

class Text : public Book {         // derived class / subclass
    string topic;
  public:
    Text(const string &t, const string &a, int n, const string &topic);
};

class Comic : public Book {
    string hero;
  public:
    Comic(const string &t, const string &a, int n, const string &hero);
};
\end{lstlisting}

\begin{center}
\begin{tikzpicture}[cls/.style={draw, minimum width=20mm, minimum height=8mm}]
\node[cls] (b) {Book};
\node[cls, below left=15mm and 5mm of b] (t) {Text};
\node[cls, below right=15mm and 5mm of b] (c) {Comic};
\draw[-{Triangle[open,length=3.5mm,width=3mm]}] (t.north) -- (b.south west);
\draw[-{Triangle[open,length=3.5mm,width=3mm]}] (c.north) -- (b.south east);
\end{tikzpicture}
\end{center}

A derived class inherits all the members (fields and methods) of its base
class. Any method you can call on a Book, you can call on a Text or Comic.

\begin{aside}
We've seen this already: anything you can do with an \code{istream} you can do
with an \code{ifstream} or \code{istringstream}, because they inherit from
\code{istream}.
\end{aside}

\subsection{Inheriting Private Members}
\code{Text} has inherited \code{title}, \code{author} and \code{numPages}, but
they are private in \code{Book}:
\begin{lstlisting}[style=cpp]
Text t{...};
t.author = "Bob";    // error: author is private
\end{lstlisting}
Private really does mean private to the class --- even subclasses can't access
it. So this constructor does not compile:
\begin{lstlisting}[style=cpp]
Text::Text(const string &t, const string &a, int n, const string &topic)
    : title{t}, author{a}, numPages{n}, topic{topic} {}   // error
\end{lstlisting}
for two reasons:
\begin{enumerate}
  \item \code{title}, \code{author}, \code{numPages} are private in Book;
  \item the MIL can only name fields declared in \emph{this} class.
\end{enumerate}

\begin{important}[Object construction, with inheritance]
\begin{enumerate}
  \item Space is allocated.
  \item The superclass part is constructed.
  \item Fields are initialized (MIL or default ctors).
  \item The ctor body runs.
\end{enumerate}
\end{important}
In step 2 the Book part is default-constructed unless told otherwise --- but Book
has no default ctor. Options: (1) give Book a default ctor; (2) call a
non-default Book ctor in the MIL:
\begin{lstlisting}[style=cpp]
Text::Text(const string &t, const string &a, int n, const string &topic)
    : Book{t, a, n}, topic{topic} {}
\end{lstlisting}
Option 2 is the right one: the superclass ctor might do work (like enforcing
invariants) that subclasses shouldn't have to repeat.

\section{Protected Visibility}
A \code{protected} member is accessible by the class and all of its subclasses
(but not by clients).
\begin{lstlisting}[style=cpp]
class Book {
  protected:
    string title, author;
    int numPages;
  public:
    ...
};
class Text : public Book {
  public:
    void addAuthor(const string &a) { author += a; }   // OK: protected
};
Text t{...};
t.author = "Bob";    // still an error for clients
\end{lstlisting}
But protected fields break encapsulation: anyone can write a subclass and
thereby get access, and a subclass may break the base class's invariants.

\begin{important}[Advice]
Keep fields private; give subclasses protected accessors/mutators instead.
\end{important}
\begin{lstlisting}[style=cpp]
class Book {
    string title, author;
    int numPages;
  protected:
    void addAuthor(const string &a) { author += a; }   // controlled access
  public:
    ...
};
\end{lstlisting}

\section*{Summary}
\begin{itemize}
  \item Design first: classes and relationships, drawn in UML.
  \item Composition (owns-a, filled diamond): parts die and copy with the whole.
  \item Aggregation (has-a, hollow diamond): parts live independently.
  \item Inheritance (is-a, open triangle): subclasses inherit all members, but
        can't touch private ones; construct the base part through the MIL.
  \item \code{protected} gives subclasses access; prefer private fields with
        protected accessors.
\end{itemize}
